% Generated draft: requires mathematical and source review.
\documentclass[11pt]{article}
\usepackage[margin=25mm]{geometry}
\usepackage{fontspec}
\setmainfont{Latin Modern Roman}
\usepackage{amsmath,amssymb,mathtools,unicode-math}
\setmathfont{Latin Modern Math}
\usepackage{enumitem,needspace,xurl,hyperref}
\setlength{\parindent}{0pt}
\setlength{\parskip}{5pt}
\setlength{\emergencystretch}{3em}
\newcommand{\N}{\mathbb{N}}
\newcommand{\Z}{\mathbb{Z}}
\newcommand{\Q}{\mathbb{Q}}
\newcommand{\R}{\mathbb{R}}
\newcommand{\C}{\mathbb{C}}
\newcommand{\F}{\mathbb{F}}
\newcommand{\A}{\mathbb{A}}
\newcommand{\PP}{\mathbb{P}}
\newcommand{\E}{\mathbb{E}}
\newcommand{\eps}{\varepsilon}
\newcommand{\dd}{\,\mathrm{d}}
\begin{document}
\section*{UnsolvedMath ID: 30004452}
\textbf{Source problem code:} catalogue entry
\subsection{Definitions and mathematical statement}
Let \(G\) be a locally compact group, with the standard convention that locally compact groups are Hausdorff, and fix a left Haar measure \(\dd s\). No discreteness or countability hypothesis is supplied. The group \(G\) is amenable if there exists a positive linear functional \(m:L^\infty(G)\to\C\) such that \(m(1)=1\) and \(m(L_s f)=m(f)\) for every \(s\in G\) and \(f\in L^\infty(G)\), where \((L_s f)(t)=f(s^{-1}t)\). It is nonamenable if no such functional exists.

Equip \(L^1(G)\) with its usual convolution and involution. For a strongly continuous unitary representation \(\pi\) of \(G\), write \(\pi(f)=\int_G f(s)\pi(s)\,\dd s\). Define the full and reduced group norms by
\[
\|f\|_{\mathrm{u}}=\sup_{\pi}\|\pi(f)\|,
\qquad
\|f\|_{\mathrm{r}}=\|\lambda_G(f)\|,
\]
where the supremum ranges over strongly continuous unitary representations and \(\lambda_G\) is the left regular representation on \(L^2(G)\), given by \((\lambda_G(s)\xi)(t)=\xi(s^{-1}t)\). Their completions are the full group algebra \(C^*(G)\) and the reduced group algebra \(C^*_{\mathrm r}(G)\), respectively.

For this draft, an \emph{exotic group \(C^*\)-algebra} means a completion \(C^*_{\mu}(G)\) of \(L^1(G)\) for a \(C^*\)-norm satisfying
\[
\|f\|_{\mathrm r}\leq\|f\|_{\mu}\leq\|f\|_{\mathrm u}
\qquad(f\in L^1(G)),
\]
with \(\|\cdot\|_{\mu}\neq\|\cdot\|_{\mathrm r}\) and \(\|\cdot\|_{\mu}\neq\|\cdot\|_{\mathrm u}\). Equivalently, the canonical surjective homomorphisms
\[
C^*(G)\longrightarrow C^*_{\mu}(G)\longrightarrow C^*_{\mathrm r}(G)
\]
are both noninjective. The two norm inequalities need not be strict on every element; each pair of norms must differ on some element.

\textbf{Question.} Does every nonamenable locally compact group \(G\) admit such a norm \(\|\cdot\|_{\mu}\), and hence an exotic group \(C^*\)-algebra?

\textbf{Unresolved source qualification.} The supplied record does not define ``exotic'' or specify additional requirements on the completion. The intermediate-completion interpretation above, including the absence of additional countability or functoriality requirements, requires confirmation against the cited report.
\subsection{Short English statement}
Does every nonamenable locally compact group have a group \(C^*\)-completion strictly intermediate between its full and reduced completions, as measured by the canonical quotient maps? This uses the intermediate-completion meaning of ``exotic''; the supplied record does not specify its definition or further hypotheses.
\subsection{Sources}
\begin{itemize}
\item Supplied local UnsolvedMath registry record, permanent ID \(30004452\), source problem code \texttt{catalogue entry}, in the Oberwolfach Reports Open Problems collection. The supplied proposed year is 2020; no individual proposer is named.
\item Supplied report citation: \emph{Geometric Structures in Group Theory} (2021), Oberwolfach Reports, pp. 877--918, DOI \texttt{10.4171/owr/2020/16}. \url{https://doi.org/10.4171/owr/2020/16}. The record describes Oberwolfach Reports as documenting mathematical workshops and their open-problem discussions, and also labels this source ``Oberwolfach Report 16/2020.'' No problem-specific page is supplied.
\item Supplied literature source: De Laat--Siebenand (2021). \url{https://doi.org/10.5802/aif.3441}. No article title or precise supporting location is supplied.
\item Supplied literature assessment, checked on exactly \texttt{2026-08-22}: the record labels the problem open and states that the locally compact question remains unresolved, that existing work develops many exotic completions, and that even the assertion for every nonamenable discrete group is explicitly left open. This is supplied dated curation derived from the OWR record, not an independently verified literature assessment. The record itself requests independent verification of time-sensitive status and cited claims.
\item The supplied statement status is \texttt{corrected\_verified}, with the explanation that mathematical content was preserved while notation, typography, omitted definitions, or duplicate clauses were repaired for self-contained publication. This is a registry qualification, not independent verification of this draft or its mathematical correctness.
\end{itemize}
\end{document}
